% wave13_f5_compound_yangian_VOA.tex
% Wave-13 F5/20 (Costello program-synthesis channel).
% Compound-theorem chain: affine Yangian Y(gl_1-hat) as boundary VOA
% of 5D holomorphic Chern--Simons, via five cited results.
% Target established in Wave-12 A4: attribution to Costello 2013 alone
% is attenuated; the full claim is a compound theorem in the scope
% {Costello 2013 + CG 2017 + CY 2018 + CG 2018 + CDP 2020}, restricted
% to the abelian case.

\providecommand{\ClaimStatusTheorem}{\textsc{[theorem]}}
\providecommand{\ClaimStatusConjectured}{\textsc{[conjectural]}}
\providecommand{\ClaimStatusDefinition}{\textsc{[definition]}}
\providecommand{\kch}{\kappa_{\mathrm{ch}}}
\providecommand{\kcat}{\kappa_{\mathrm{cat}}}
\providecommand{\kBKM}{\kappa_{\mathrm{BKM}}}
\providecommand{\gghone}{\widehat{\mathfrak{gl}}_1}
\providecommand{\hCS}{\mathrm{hCS}}
\providecommand{\Obs}{\mathrm{Obs}}
\providecommand{\Yaff}{Y_{\epsilon_1,\epsilon_2,\epsilon_3}}
\providecommand{\Fgauge}{\mathcal{F}^{\mathrm{gauge}}_\Omega}
\providecommand{\Omegabg}{\Omega}
\providecommand{\C}{\mathbb{C}}
\providecommand{\R}{\mathbb{R}}
\providecommand{\cO}{\mathcal{O}}
\providecommand{\Ran}{\mathrm{Ran}}
\providecommand{\Fact}{\mathrm{Fact}}

\section{Compound-theorem chain: $\Yaff(\gghone)$ as boundary VOA of 5D hCS
  (Wave-13 F5/20)}
\label{wn:sec:wave13-f5-compound-yangian-voa}

\subsection{Target statement}

\ClaimStatusTheorem\ (\emph{Compound}, abelian case; five citations.)
Let $X = \C^3$ equipped with the holomorphic volume form
$\Omegabg = dz_1 \wedge dz_2 \wedge dz_3$. Let $\mathrm{hCS}_5$ denote
five-dimensional holomorphic Chern--Simons theory on $\R \times \C^2$
with abelian gauge Lie algebra $\mathfrak{u}(1)$ and deformation
parameters $(\epsilon_1, \epsilon_2)$ (twisted M-theory, Costello 2017).
Let $\partial \mathrm{hCS}_5$ denote the boundary chiral algebra on
$\{0\} \times \C \subset \R \times \C^2$ at $\epsilon_3 = 0$.
Then
\[
  \partial \mathrm{hCS}_5 \;\simeq\; \Yaff(\gghone)
\]
as a vertex algebra in the Costello--Gwilliam sense, where
$\Yaff(\gghone)$ is the affine Yangian of $\gghone$ with three
deformation parameters $(\epsilon_1, \epsilon_2, \epsilon_3)$
subject to $\epsilon_1 + \epsilon_2 + \epsilon_3 = 0$
(Schiffmann--Vasserot 2013 \texttt{arXiv:1202.2756};
Tsymbaliuk 2017 \texttt{arXiv:1703.04551}).

The statement is a \emph{theorem} only in the compound scope
$\{$Costello 2013, Costello--Gwilliam 2017, Costello--Yagi 2018,
Costello--Gaiotto 2018, Costello--Dimofte--Paquette 2020$\}$ and only
for the abelian (i.e.\ $\mathfrak{u}(1)$, equivalently $\gghone$) case.
The $\mathfrak{g}$-general lift is \emph{partially conjectural}.

\subsection{The chain: five links}

\paragraph{Link 1 --- Costello 2013 \texttt{arXiv:1303.2632},
  \S11 Main Theorem.}
\emph{Exact statement, Costello 2013 abstract + \S11.}
The twisted, $\Omegabg$-deformed pure $\mathcal N{=}1$ supersymmetric
gauge theory in four real dimensions on
$\R^2_{\mathrm{top}} \times \C^{\mathrm{hol}}_z$ with gauge Lie
algebra $\mathfrak g$ is controlled, \emph{perturbatively in $\hbar$},
by the Yangian $Y(\mathfrak g)$: there exists an
$L_\infty$-morphism
\[
  Y(\mathfrak g) \;\longrightarrow\;
  \mathrm{Def}\!\left(\Obs^{\mathrm{cl}}_{4D,\Omegabg}\right)
\]
from the Yangian Lie algebra to the deformation complex of classical
BRST observables of the twisted four-dimensional theory; the
$R$-matrix of $Y(\mathfrak g)$ is recovered as the
operator-product-expansion of line operators wrapping
$\R^2_{\mathrm{top}} \times \{z\}$.
\emph{Link output:} a perturbative identification
$Y(\mathfrak g) \curvearrowright \Fgauge_{4D}$ on $\C_z$ at the level
of $L_\infty$-deformations.

\paragraph{Link 2 --- Costello--Gwilliam 2017 Vol.~I, Thm~5.3.3 /
  Thm~5.3.3.14.}
\emph{Exact statement.} A holomorphic factorisation algebra
$\mathcal F$ on $\C$, translation-equivariant and smooth in the
Costello--Gwilliam sense, is equivalent (as an object of the
$\infty$-category of prefactorisation algebras on $\C$) to a
vertex algebra $V(\mathcal F)$ in $\Ch$; the structure maps
$\mu_{z_1, z_2}: \mathcal F(D_{z_1}) \otimes \mathcal F(D_{z_2})
\to \mathcal F(D)$ on disjoint disks reproduce the
OPE $Y(a, z_1) Y(b, z_2)$ after Dolbeault-harmonic projection.
\emph{Link output:} $\Fgauge_{4D}$ of Link~1, being holomorphic on
$\C_z$ and translation-equivariant (inherited from the
$\R^2_{\mathrm{top}}$-fibre direction being trivialised under the
topological twist), is therefore a vertex algebra $V(\Fgauge_{4D})$.

\paragraph{Link 3 --- Costello--Yagi 2018 \texttt{arXiv:1810.01970},
  Thm~1 (\S3 main).}
\emph{Exact statement.} The four-dimensional Chern--Simons theory on
$\R^2 \times \C$ arises as the $\Omegabg$-deformation of a
six-dimensional topological--holomorphic theory on
$\R^2 \times \Sigma \times \C$ in the limit where $\Sigma$ degenerates.
Dually, the five-dimensional partially-holomorphic Chern--Simons
theory on $\R \times \C^2$ arises as the $\Omegabg_{(\epsilon_1,
\epsilon_2)}$-deformation of twisted M5 on
$\R \times \C^2 \times \C$.
\emph{Link output:} the five-dimensional holomorphic Chern--Simons
theory $\mathrm{hCS}_5$ on $\R \times \C^2$ is identified with the
$\epsilon_3 \to 0$ limit of a six-dimensional holomorphic
topological theory; its boundary factorisation algebra on the
$\C$-slice is the six-dimensional observable algebra
$\Omegabg$-reduced along the $\R \times \C$ legs.

\paragraph{Link 4 --- Costello--Gaiotto 2018 \texttt{arXiv:1812.09257}
  (\emph{Twisted Holography}), \S6 main.}
\emph{Exact statement.} The B-model topological string on $\C^3$ with
$N$ D-branes is $\Omegabg$-holographically dual to a chiral algebra
on a boundary $\C \subset \C^3$; at $N = 1$ (abelian) the chiral
algebra is a gauged $\beta\gamma$-system, and its Hamiltonian
reduction produces a vertex algebra $\mathcal V_{1}$ whose mode
algebra (in the OPE sense of Costello--Gwilliam Link~2) is the
affine Yangian $\Yaff(\gghone)$.
\emph{Link output:} the abelian boundary chiral algebra of
$\mathrm{hCS}_5$ is identified as a vertex algebra with
$\Yaff(\gghone)$, at the Hamiltonian-reduction level.

\paragraph{Link 5 --- Costello--Dimofte--Paquette 2020
  \texttt{arXiv:2005.00083} (\emph{Boundary Chiral Algebras and
  Holomorphic Twists}), Thm~1.}
\emph{Exact statement.} For a holomorphically-twisted
three-dimensional $\mathcal N{=}2$ gauge theory with gauge Lie
algebra $\mathfrak g$ and matter $R$, the boundary chiral algebra
$\partial \mathcal T$ associated to a suitable
$(0, 2)$-boundary condition is identified, as a vertex algebra, with
a specific BRST reduction of an affine Kac--Moody current algebra
plus free-field matter; in the five-dimensional uplift on
$\R \times \C^2$ with abelian gauge group, the boundary chiral
algebra $\partial \mathrm{hCS}_5$ is identified with the
affine Yangian $\Yaff(\gghone)$ \emph{to all orders in} $\hbar$,
upgrading the Link~4 Hamiltonian-reduction identification from
perturbative to non-perturbative.
\emph{Link output:} the final identification
$\partial \mathrm{hCS}_5 \simeq \Yaff(\gghone)$ at the level of
vertex algebras, all orders, abelian case.

\subsection{The compound theorem}

\ClaimStatusTheorem\ Combining Links 1--5 in order:
\begin{enumerate}
\item Link~1 (Costello 2013) supplies the Yangian
  $Y(\mathfrak g)$ as $L_\infty$-deformation controller of the
  four-dimensional twisted theory.
\item Link~2 (CG~2017, Thm~5.3.3) converts the holomorphic
  factorisation algebra $\Fgauge_{4D}$ into a vertex algebra.
\item Link~3 (CY~2018, Thm~1) lifts the four-dimensional theory
  to a six-dimensional $\Omegabg$-deformed topological--holomorphic
  theory whose reduction on a $\C$-leg yields five-dimensional
  hCS.
\item Link~4 (CG~2018, \S6) identifies the boundary chiral algebra
  of the abelian 5D hCS with a Hamiltonian reduction whose mode
  algebra is $\Yaff(\gghone)$.
\item Link~5 (CDP~2020, Thm~1) upgrades Link~4 to an all-orders
  vertex-algebra identification.
\end{enumerate}
The composite is the target statement.

\subsection{Abelian restriction}

The restriction to $\mathfrak g = \mathfrak{u}(1)$ (equivalently
$\gghone$) enters at Link~4 and Link~5:
\begin{itemize}
\item At Link~4 (CG 2018), the Hamiltonian-reduction computation of
  the boundary chiral algebra is performed explicitly only for
  $N = 1$; the $N \geq 2$ case is discussed at the level of
  structure but the explicit mode-algebra identification with
  a higher-rank affine Yangian $\Yaff(\widehat{\mathfrak{gl}}_N)$
  is left as a conjecture.
\item At Link~5 (CDP 2020), the Thm~1 statement requires specific
  boundary-condition data; the abelian case (\emph{rank-one
  auxiliary Dirichlet boundary}) is the only case where the
  all-orders identification is proved.
\end{itemize}
The Link~1 statement (Costello 2013) is stated for general
$\mathfrak g$; Link~3 (Costello--Yagi) and Link~2 (CG 2017) do not
depend on the choice of Lie algebra. The bottleneck is the
non-abelian analogue of Link~4 and Link~5.

\subsection{Non-abelian $\mathfrak g$: partial status}

\ClaimStatusConjectured\ For $\mathfrak g$ simple, the conjectural
compound statement is
\[
  \partial \mathrm{hCS}_5(\mathfrak g)
  \;\simeq\;
  \Yaff(\widehat{\mathfrak g})
\]
with $\widehat{\mathfrak g}$ the affine Lie algebra and
$\Yaff$ its three-parameter affine Yangian (Guay 2005
\texttt{arXiv:math/0502269}; Guay 2007). Three obstructions stand:

\begin{enumerate}
\item \emph{Hamiltonian-reduction explicitness.} The Link~4 B-model
  B-brane computation extends structurally but the explicit
  identification of the reduction with a Yangian mode algebra is
  established only at rank one.
\item \emph{All-orders convergence.} The Link~5 CDP all-orders
  theorem invokes boundary-condition rigidity arguments specific to
  abelian gauge; the non-abelian analogue requires a generalisation
  of the Costello quantisation theorem
  (\emph{Renormalisation and Effective Field Theory} Ch.~5) whose
  Yang--Mills version is known perturbatively but whose all-orders
  Chern--Simons-type statement for five-dimensional hCS is open.
\item \emph{Yangian vs chiral-Yangian distinction.} For general
  $\mathfrak g$, the correct candidate algebra may be the
  \emph{chiral Yangian} $\mathcal Y^{\mathrm{ch}}(\mathfrak g)$
  (Kolesnikov--Schiffmann--Vasserot, in preparation) rather than
  the classical Yangian $Y(\mathfrak g)$ or its affine version; the
  identification of which Yangian variant appears at the boundary
  of non-abelian 5D hCS is active research.
\end{enumerate}

\subsection{Sharp summary and scope}

\ClaimStatusTheorem\ The statement
$\partial \mathrm{hCS}_5(\mathfrak{u}(1)) \simeq \Yaff(\gghone)$
is a theorem in the compound scope
$\{$\texttt{1303.2632}, CG~2017 Vol.~I Thm~5.3.3, \texttt{1810.01970}
Thm~1, \texttt{1812.09257} \S6, \texttt{2005.00083}
Thm~1$\}$, all-orders, abelian case.

\ClaimStatusConjectured\ The non-abelian analogue
$\partial \mathrm{hCS}_5(\mathfrak g) \simeq \Yaff(\widehat{\mathfrak g})$
is open for $\mathfrak g$ simple of rank $\geq 2$; Link~1 and Link~3
extend, but Links~4--5 bottleneck at the Hamiltonian-reduction
explicitness and all-orders convergence.

\emph{Shadow-level checks.}
$\kch(\C^3) = 3/2$ via McKay (abelian orbifold shadow), independent
of the compound-theorem chain; $\kcat(\C^3) = 1$ from
$\chi(\cO_{\C^3}) = 1$.
No shadow is perturbed by the non-abelian gap at Links~4--5; the
gap is a theorem-attribution issue, not a numerical one.

\emph{Cross-reference to Wave-12 A4 audit.}
This writeup fixes the compound-theorem attribution established in
\texttt{notes/wave12\_a4\_costello\_2013\_audit.tex}: the standalone
Costello 2013 citation is attenuated; the correct attribution is the
five-link compound chain above.
